% ============================================================================
%  Documento de costos frente a venta. Informe 3.
%  Formulario T-22 (FEP01 p.69): proveedores clave, adquisiciones clave,
%  curva S, análisis de costos, VAN y TIR, y valorización de las cinco
%  innovaciones en el flujo de caja. Las cifras salen de
%  economico/datos/partidas.csv: aquí no se escribe ningún monto a mano.
% ============================================================================

\begin{resumenApertura}
\marcador{Qué muestra este documento, en cuatro o cinco líneas, con el valor
ofertado frente al costo total: \textbackslash montoPartida\{proyecto\} y
\textbackslash montoPartida\{gasto-total\}}
\begin{recibe}
  \item \marcador{Qué recibe el mandante}
\end{recibe}
\end{resumenApertura}

\section{Proveedores clave de la solución}\label{sec:eco-proveedores}

\marcador{Párrafo de contexto: qué muestra la \tabref{eco-proveedores} y por qué
esos proveedores son clave}

\begin{tabla}{Proveedores clave de la solución}{eco-proveedores}{L{38mm} X L{38mm} L{34mm}}
Proveedor & Qué provee & Componente o servicio & Criterio de selección \\ \midrule
\marcador{proveedor} & \marcador{producto o servicio} & \marcador{componente} & \marcador{criterio} \\
\end{tabla}

\marcador{Conclusión que se saca de la tabla}

\section{Adquisiciones clave o relevantes}\label{sec:eco-adquisiciones}

\marcador{Párrafo de contexto: qué muestra la \tabref{eco-adquisiciones}. El
monto de cada adquisición va en partidas.csv y se imprime con
\textbackslash montoPartida}

\begin{tabla}{Adquisiciones clave o relevantes de la solución}{eco-adquisiciones}{L{40mm} X C{14mm} L{36mm}}
Adquisición & Para qué & Mes & Modalidad \\ \midrule
\marcador{adquisición} & \marcador{uso en la solución} & \marcador{mes} & \marcador{compra, arriendo o suscripción} \\
\end{tabla}

\marcador{Conclusión que se saca de la tabla}

\section{Curva S del proyecto}\label{sec:eco-curva-s}

\marcador{Párrafo de contexto: costos e ingresos acumulados por mes, del mes 1
al 56, horizonte de evaluación del Formulario E-24. La figura sale de la
planilla y se cita con \textbackslash figref antes de aparecer}

\section{Análisis de costos de la solución}\label{sec:eco-costos}

\marcador{Párrafo de contexto: qué muestra la \tabref{eco-gastos}. El Formulario
T-22 pide los gastos desglosados en operación, inversión, administrativos y
recursos humanos, con los gastos de proveedores sumados a operación}

\tablaMontos{Costos de la solución por categoría de gasto}{eco-gastos}%
  {gasto-operacion-directo, gasto-proveedores, gasto-inversion, gasto-administracion, gasto-rrhh}

\marcador{Justificación de cada ítem y su coherencia con el plan de
actividades: las bases evalúan con severidad un presupuesto sobreestimado o
subestimado}

\subsection{Costos frente a venta}\label{sec:eco-venta}

\marcador{Párrafo de contexto: qué muestra la \tabref{eco-venta}, el valor
ofertado que se compara con el costo total de la \tabref{eco-gastos}}

\tablaMontos{Valor ofertado por fase}{eco-venta}{implementacion, operacion}

\marcador{Margen resultante y de dónde sale}

\section{VAN y TIR de la solución}\label{sec:eco-van}

\marcador{Párrafo de contexto: qué muestra la \tabref{eco-van} y con qué tasa de
descuento se calcula el VAN}

\begin{tabla}{Indicadores de la evaluación financiera}{eco-van}{L{60mm} X}
Indicador & Valor \\ \midrule
Horizonte de evaluación & 56 meses, \form{E-24}{73} \\
Tasa de descuento & \marcador{tasa y justificación} \\
Valor actual neto (VAN) & \montoPartida{van} \\
Tasa interna de retorno (TIR) & \marcador{porcentaje} \\
\end{tabla}

\marcador{Conclusión que se saca de los indicadores}

\section{Valorización de las cinco innovaciones}\label{sec:eco-innovaciones}

\marcador{Párrafo de contexto: qué muestra la \tabref{eco-innovaciones}, la
inversión, el costo operacional y el beneficio esperado de cada innovación en
el flujo de caja}

\tablaMontos*[fuente=condensada]{Valorización de las cinco innovaciones en el flujo de caja}{eco-innovaciones}%
  {innovacion-1-inversion, innovacion-1-costo, innovacion-1-beneficio,
   innovacion-2-inversion, innovacion-2-costo, innovacion-2-beneficio,
   innovacion-3-inversion, innovacion-3-costo, innovacion-3-beneficio,
   innovacion-4-inversion, innovacion-4-costo, innovacion-4-beneficio,
   innovacion-5-inversion, innovacion-5-costo, innovacion-5-beneficio}

\marcador{Conclusión que se saca de la tabla}

\section{Planilla de cálculo}\label{sec:eco-planilla}

\marcador{Qué contiene la planilla que entrega el mandante: gastos de operación,
inversión, administrativos y recursos humanos, una hoja con los gastos de
proveedores sumados a operación y el flujo de caja mensual con total del mes y
monto acumulado. Nombre del archivo y hoja donde está cada valor de este
documento}
